\section{Nested square patches}
\label{sec:peps_nested_patch_geometry}

Fix a finite induced lattice domain $\Lambda\subset\mathbb Z^2$ and a centre
$c\in\mathbb R^2$. The square patches in the proof of
\cite[Proposition~4.1]{OpenAI2026PolynomialPEPS} are samples of closed ambient
squares. Intersecting with $\Lambda$ permits clipping at the physical boundary
and does not add edges to the induced graph.
% This section covers the geometric part of 03-patches.tex only. It does not
% mark Proposition 4.1, its energy estimate, or its rank conclusion as proved.

\begin{definition}[Closed square sample]
    \label{def:peps_closed_square_sample}
    \lean{
        TNLean.PEPS.AreaLaw.Geometry.closedSquareSample,
        TNLean.PEPS.AreaLaw.Geometry.mem_closedSquareSample,
        TNLean.PEPS.AreaLaw.Geometry.mem_closedSquareSample_iff_int,
        TNLean.PEPS.AreaLaw.Geometry.closedSquareSample_mono
    }
    \leanok
    \uses{def:area_law_domain}
    For $r\in\mathbb R$, define
    \begin{align}
        K(c,r) & =\Lambda\cap
        \bigl([c_1-r,c_1+r]\times[c_2-r,c_2+r]\bigr).
        \label{eq:peps_closed_square_sample}
    \end{align}
    Thus $x\in K(c,r)$ exactly when $x\in\Lambda$ and
    $|x_1-c_1|\le r$, $|x_2-c_2|\le r$. Sites on the square boundary are
    included. Equivalently, each integer coordinate lies between
    $\lceil c_a-r\rceil$ and $\lfloor c_a+r\rfloor$, for $a=1,2$.
    The sample is empty for $r<0$ and increases with $r$.
\end{definition}

\begin{definition}[Distance from the square centre]
    \label{def:peps_square_radius}
    \lean{
        TNLean.PEPS.AreaLaw.Geometry.squareRadius,
        TNLean.PEPS.AreaLaw.Geometry.mem_closedSquareSample_iff_radius
    }
    \leanok
    \uses{def:peps_closed_square_sample}
    For a lattice point $x$, set
    $\rho(x)=\max\{|x_1-c_1|,|x_2-c_2|\}$.
    Then $K(c,r)=\{x\in\Lambda:\rho(x)\le r\}$.
\end{definition}

\begin{theorem}[Nearest neighbours and closed-square crossing]
    \label{thm:peps_square_crossing_interval}
    \lean{
        TNLean.PEPS.AreaLaw.Geometry.abs_coordinate_sub_le_one_of_adj,
        TNLean.PEPS.AreaLaw.Geometry.mem_closedSquareSample_of_adj,
        TNLean.PEPS.AreaLaw.Geometry.abs_squareRadius_sub_le_one_of_adj,
        TNLean.PEPS.AreaLaw.Geometry.mem_edgeBoundary_pair_iff,
        TNLean.PEPS.AreaLaw.Geometry.mem_edgeBoundary_closedSquareSample_iff
    }
    \leanok
    \uses{def:peps_square_radius,def:area_law_hamiltonian}
    Adjacent sites $x,y$ satisfy $|x_a-y_a|\le1$ for $a=1,2$ and
    $|\rho(x)-\rho(y)|\le1$. In particular, if $x\in K(c,r)$, then
    $y\in K(c,r+1)$. For any cut $A$, an unordered pair belongs to
    $\partial_\Lambda A$ precisely when its sites are adjacent and exactly
    one belongs to $A$. Consequently, for adjacent $x,y$,
    \begin{align}
        \{x,y\}\in\partial_\Lambda K(c,r)
         & \quad\Longleftrightarrow\quad
        \min\{\rho(x),\rho(y)\}\le r<\max\{\rho(x),\rho(y)\}.
        \label{eq:peps_square_crossing_interval}
    \end{align}
    The interval is closed at the smaller radial distance and open at the
    larger radial distance.
\end{theorem}
\begin{proof}\leanok
    One coordinate of adjacent sites differs by one and the other is equal,
    so $|x_a-y_a|\le1$. For $a=1,2$ the triangle inequality gives
    \begin{align*}
        |y_a-c_a|\le|x_a-c_a|+|x_a-y_a|\le\rho(x)+1,
    \end{align*}
    hence $\rho(y)\le\rho(x)+1$; interchanging $x$ and $y$ gives
    $|\rho(x)-\rho(y)|\le1$. An unordered pair crosses $A$ exactly when
    \begin{align*}
        (x\in A\wedge y\notin A)\vee(y\in A\wedge x\notin A).
    \end{align*}
    For $A=K(c,r)$, membership of a site $z$ is $\rho(z)\le r$, and the
    disjunction becomes
    $(\rho(x)\le r<\rho(y))\vee(\rho(y)\le r<\rho(x))$, which is
    (\ref{eq:peps_square_crossing_interval}).
\end{proof}

\begin{theorem}[Number of sites in a sampled square]
    \label{thm:peps_closed_square_card}
    \lean{
        TNLean.PEPS.AreaLaw.Geometry.card_closedSquareSample_le,
        TNLean.PEPS.AreaLaw.Geometry.card_closedSquareSample_le_eightyOne
    }
    \leanok
    \uses{def:peps_closed_square_sample}
    If $r\ge0$, then $|K(c,r)|\le(2r+1)^2$.
    If also $r\le4$, then $|K(c,r)|\le81$.
\end{theorem}
\begin{proof}\leanok
    Each coordinate interval has at most $2r+1$ integer points. The sample
    is a subset of their Cartesian product. When $r\le4$, its cardinality is
    at most $9^2=81$.
\end{proof}

\begin{theorem}[Crossing edges of a sampled square]
    \label{thm:peps_closed_square_boundary}
    \lean{TNLean.PEPS.AreaLaw.Geometry.card_edgeBoundary_closedSquareSample_le}
    \leanok
    \uses{def:peps_closed_square_sample,def:area_law_hamiltonian}
    For every $r\ge0$, the number of unordered induced edges crossing the
    sampled square satisfies
    $|\partial_\Lambda K(c,r)|\le8r+4$.
\end{theorem}
\begin{proof}\leanok
    A horizontal crossing edge lies on an integer row meeting
    $[c_2-r,c_2+r]$. On each such row there are at most two crossing edges,
    one at each end of the integer interval inside the square. An interval of
    length $2r$ contains at most $2r+1$ integers. The same argument applies to
    vertical edges. Hence the total is at most
    $2(2r+1)+2(2r+1)=8r+4$. Removing sites outside $\Lambda$ cannot create
    an induced crossing edge.
\end{proof}

\begin{definition}[Nested square patches]
    \label{def:peps_nested_patches}
    \lean{
        TNLean.PEPS.AreaLaw.Geometry.nestedPatchCount,
        TNLean.PEPS.AreaLaw.Geometry.nestedPatchRadius,
        TNLean.PEPS.AreaLaw.Geometry.nestedPatch
    }
    \leanok
    \uses{def:peps_closed_square_sample}
    For $u\ge0$, put $m=\lfloor u/2\rfloor+1$. For each integer
    $0\le j<m$, define
    \begin{align}
        r_j & =u+2j,\qquad K_j=K(c,r_j).
        \label{eq:peps_nested_patches}
    \end{align}
\end{definition}

\begin{theorem}[Number, radii, and boundary size of the patches]
    \label{thm:peps_nested_patch_sizes}
    \lean{
        TNLean.PEPS.AreaLaw.Geometry.nestedPatchCount_pos,
        TNLean.PEPS.AreaLaw.Geometry.half_lt_nestedPatchCount,
        TNLean.PEPS.AreaLaw.Geometry.le_nestedPatchRadius,
        TNLean.PEPS.AreaLaw.Geometry.nestedPatchRadius_le_two_mul,
        TNLean.PEPS.AreaLaw.Geometry.nestedPatchRadius_add_two_le,
        TNLean.PEPS.AreaLaw.Geometry.nestedPatch_mono,
        TNLean.PEPS.AreaLaw.Geometry.card_edgeBoundary_nestedPatch_le
    }
    \leanok
    \uses{def:peps_nested_patches,def:area_law_hamiltonian}
    The family contains $m\ge1$ patches and $m>u/2$. Every selected radius
    satisfies $u\le r_j\le2u$ and every selected patch satisfies
    $|\partial_\Lambda K_j|\le16u+4$.
    For indices $i<j$, one has $r_i+2\le r_j$. For $i\le j$, one has
    $K_i\subseteq K_j$; clipping can make these inclusions equalities.
\end{theorem}
\begin{proof}\leanok
    \uses{thm:peps_closed_square_boundary}
    The floor inequalities give $u/2<m$ and $j\le\lfloor u/2\rfloor\le u/2$.
    Thus $u\le u+2j\le2u$, and the crossing-edge estimate applies.
    Increasing an integer index by at least one increases the radius by at
    least two; inclusion follows from the closed coordinate inequalities.
\end{proof}

\begin{theorem}[Supports of crossing edges under radius separation]
    \label{thm:peps_square_crossing_support}
    \lean{
        TNLean.PEPS.AreaLaw.Geometry.edgeBoundary_support_disjoint_of_add_one_le,
        TNLean.PEPS.AreaLaw.Geometry.edgeBoundary_support_subset_of_add_one_le
    }
    \leanok
    \uses{def:peps_closed_square_sample,def:area_law_hamiltonian}
    Suppose $e\in\partial_\Lambda K(c,r)$. If $s+1\le r$, then the two
    sites of $e$ are disjoint from $K(c,s)$. If $r+1\le t$, then both sites
    belong to $K(c,t)$.
\end{theorem}
\begin{proof}\leanok
    \uses{thm:peps_square_crossing_interval}
    Write $e=\{x,y\}$ with $x\in K(c,r)$ and $y\notin K(c,r)$.
    If either site belonged to $K(c,s)$, the nearest-neighbour enlargement
    would put $y$ in $K(c,s+1)\subseteq K(c,r)$, a contradiction.
    Conversely, $x\in K(c,r)$ implies $y\in K(c,r+1)$, so both sites
    belong to every $K(c,t)$ with $r+1\le t$.
\end{proof}

\begin{theorem}[A single crossing among nested patches]
    \label{thm:peps_nested_patch_crossing}
    \lean{
        TNLean.PEPS.AreaLaw.Geometry.nestedPatch_crossing_disjoint_earlier,
        TNLean.PEPS.AreaLaw.Geometry.nestedPatch_crossing_subset_later,
        TNLean.PEPS.AreaLaw.Geometry.nestedPatch_crossing_unique,
        TNLean.PEPS.AreaLaw.Geometry.disjoint_edgeBoundary_nestedPatch
    }
    \leanok
    \uses{def:peps_nested_patches,def:area_law_hamiltonian}
    Let $e=\{x,y\}$ be an induced nearest-neighbour edge. If
    $e\in\partial_\Lambda K_j$, then, for every $i<j$,
    $\{x,y\}\cap K_i=\varnothing$, and, for every $j<k$,
    $\{x,y\}\subseteq K_k$. Consequently, an edge crosses at most one patch
    in the family, and the crossing-edge sets of distinct patches are disjoint.
\end{theorem}
\begin{proof}\leanok
    \uses{thm:peps_nested_patch_sizes,thm:peps_square_crossing_support}
    Earlier and later radii differ from $r_j$ by at least two, so the
    separated-radius support statements apply. If an edge crossed two
    patches, order their distinct indices. Crossing the earlier patch places
    both sites in the later patch, contradicting crossing of the later one.
\end{proof}

The following finite example uses
$\Lambda=\{0,\ldots,10\}\times\{0,1,2\}$,
$c=(5/4,3/2)$, and $u=4$. Its three clipped samples are
$K_0=\{0,\ldots,5\}\times\{0,1,2\}$,
$K_1=\{0,\ldots,7\}\times\{0,1,2\}$, and
$K_2=\{0,\ldots,9\}\times\{0,1,2\}$.
The marked edge $e=\{(7,1),(8,1)\}$ crosses $K_1$; it is disjoint from $K_0$
and contained in $K_2$. The rounded outlines enclose the clipped sets of
physical sites; they are not the boundaries of the ambient real squares.
% Formula verdict: exact samples of the three closed squares at the stated
% nonintegral centre, with radii 4, 6, 8. The finite domain clips every square.
% Ink-to-index verdict: dots denote lattice sites and lines denote adjacency;
% this is a graph illustration, not a tensor contraction or an operator.
% Contraction/boundary verdict: 33 sites, 52 distinct internal adjacency
% segments, no periodic edges, physical legs, or edges outside the domain.
% Source verdict: original finite example of the geometric step in
% September 24, 2026, 03-patches.tex, Proposition 4.1.
\begin{center}
    \tenkzkernel
    \begin{tenkz}[rows={wire,wire,wire}, cols=11, metrics=compact,
            west=none, east=none, north=none, south=none]
        \tn[skin=dot]{} & \tn[skin=dot]{} & \tn[skin=dot]{} & \tn[skin=dot]{} & \tn[skin=dot]{} & \tn[skin=dot]{} & \tn[skin=dot]{} & \tn[skin=dot]{} & \tn[skin=dot]{} & \tn[skin=dot]{} & \tn[skin=dot]{} \\
        \tn[skin=dot]{} & \tn[skin=dot]{} & \tn[skin=dot]{} & \tn[skin=dot]{} & \tn[skin=dot]{} & \tn[skin=dot]{} & \tn[skin=dot]{} & \tn[at=(2,8), name=inside, skin=dot, species=marked]{} & \tn[at=(2,9), name=outside, skin=dot, species=marked]{} & \tn[skin=dot]{} & \tn[skin=dot]{} \\
        \tn[skin=dot]{} & \tn[skin=dot]{} & \tn[skin=dot]{} & \tn[skin=dot]{} & \tn[skin=dot]{} & \tn[skin=dot]{} & \tn[skin=dot]{} & \tn[skin=dot]{} & \tn[skin=dot]{} & \tn[skin=dot]{} & \tn[skin=dot]{}
        \tnwire[name=crossing, species=marked]{inside.0}{outside.180}
        \tnmark[form=enclosure, species=selected, label pos=s]
        {(1,1) .. (3,6)}{$\textcolor{tenkzAction}{K_0}$}
        \tnmark[form=enclosure, species=secondary, label pos=n]
        {(1,1) .. (3,8)}{$\textcolor{tenkzMarked}{K_1}$}
        \tnmark[form=enclosure, species=collar, label pos=se]
        {(1,1) .. (3,10)}{$\textcolor{tenkzExtra}{K_2}$}
        \tnmark[form=label, label pos=n]{on crossing 0.65}{$e$}
    \end{tenkz}
\end{center}


\begin{theorem}[Uniform quadratic boundary sum]
    \label{thm:peps_nested_patch_quadratic_sum}
    \lean{
        TNLean.PEPS.AreaLaw.Geometry.nestedPatch_quadratic_budget,
        TNLean.PEPS.AreaLaw.Geometry.sum_card_edgeBoundary_nestedPatch_mul_sq_le
    }
    \leanok
    \uses{def:peps_nested_patches,def:area_law_hamiltonian}
    Let $u\ge4$, $\kappa\in\mathbb R$, and assign each selected patch the
    same weight $t_j=\kappa/m$. Then
    \begin{align}
        (16u+4)\sum_{j=0}^{m-1}t_j^2                & \le36\kappa^2,
        \label{eq:peps_nested_patch_uniform_budget}                  \\
        \sum_{j=0}^{m-1}|\partial_\Lambda K_j|t_j^2 & \le36\kappa^2.
        \label{eq:peps_nested_patch_quadratic_sum}
    \end{align}
\end{theorem}
\begin{proof}\leanok
    \uses{thm:peps_nested_patch_sizes}
    Since $m>u/2$ and $u\ge4$, one has $16u+4\le36m$. The constant
    weights satisfy $\sum_{j=0}^{m-1}t_j^2=\kappa^2/m$, which gives
    (\ref{eq:peps_nested_patch_uniform_budget}). Bounding each boundary
    size by $16u+4$ gives
    \begin{align*}
        \sum_{j=0}^{m-1}|\partial_\Lambda K_j|t_j^2
         & \le(16u+4)\sum_{j=0}^{m-1}t_j^2,
    \end{align*}
    and (\ref{eq:peps_nested_patch_uniform_budget}) gives
    (\ref{eq:peps_nested_patch_quadratic_sum}).
\end{proof}
% The uniform quadratic sum establishes a geometric coefficient only. It
% does not identify this sum with a Hamiltonian energy increment.
